\documentclass[11pt,a4paper]{article}

\usepackage[utf8]{inputenc}
\usepackage[T1]{fontenc}
\usepackage[margin=1in]{geometry}
\usepackage{amsmath,amssymb,amsthm}
\usepackage{mathtools}
\usepackage{booktabs}
\usepackage{hyperref}

\newtheorem{theorem}{Theorem}
\newtheorem{lemma}[theorem]{Lemma}
\newtheorem{proposition}[theorem]{Proposition}
\newtheorem{corollary}[theorem]{Corollary}
\theoremstyle{definition}
\newtheorem{definition}[theorem]{Definition}
\newtheorem{remark}[theorem]{Remark}
\theoremstyle{remark}
\newtheorem*{note}{Note}

\newcommand{\arr}{\mathsf{arr}}
\newcommand{\comp}{\mathbin{>\!\!>}}
\newcommand{\first}{\mathsf{first}}
\newcommand{\leftc}{\mathsf{left}}
\newcommand{\loopop}{\mathsf{loop}}
\newcommand{\app}{\mathsf{app}}
\newcommand{\ArrowApply}{\mathsf{ArrowApply}}
\newcommand{\fix}{\mathsf{fix}}
\newcommand{\T}{\mathcal{T}}
\newcommand{\interp}{\mathcal{I}}
\newcommand{\Sinv}{\mathcal{S}}
\newcommand{\lamterm}{\Lambda}
\newcommand{\arrowprog}{\mathcal{A}}
\newcommand{\beteta}{\equiv_{\beta\eta}}
\newcommand{\betetafix}{\equiv_{\beta\eta\fix}}
\newcommand{\stlc}{\lambda_{\times,+,\fix}}
\newcommand{\size}[1]{\lvert #1 \rvert}
\newcommand{\fanout}{\mathbin{\&\!\!\&}}
\newcommand{\lam}[2]{\lambda #1.\,#2}

\title{No Translation Asymmetry Between Arrows and Lambda:\\
A Negative Result, with a Formal Boundary for Structural Observability}

\author{}
\date{}

\begin{document}
\maketitle

\begin{abstract}
We investigate whether arrows~\cite{hughes2000} can serve as a foundational
calculus for programming languages by asking a precise question: do arrow
programs have provably awkward lambda encodings---encodings that suffer size
blowup or non-compositionality in a way that would argue for arrows as a
foundation rather than a notational convenience? We prove a \emph{negative}
result: the standard arrow-to-lambda translation $\T$ is linear
($\size{\T(p)} \leq 8\cdot\size{p}$, mechanized in Lean~4), compositional
($\T(f \comp g) = \T(g) \circ \T(f)$ definitionally), and law-preserving
($\first$ and $\leftc$ commute with composition up to $\beta\eta$, and the
loop extension law holds under $\beta\eta$ plus fix-fusion/Beki\'c's lemma).
The one subtlety is that the loop extension law requires fix-fusion, not
just $\beta\eta$; but this is a general property of recursion, not of arrows.

The asymmetries that \emph{do} exist are of two other kinds:
\emph{representation} (arrow trees are inspectable; opaque monadic closures
are not) and \emph{expressive power} (the inverse translation requires
$\ArrowApply \cong \mathsf{Monad}$). We provide three lines of evidence for
the practical significance of the representation asymmetry: (1)~experimental
(205 Rust tests): arrow representations detect independence by construction
while naive lambda analysis fails on obscured code and sophisticated lambda
recovers structure at 2.3$\times$ the cost; (2)~engineering (221 Rust
tests): speculative compilation of transparent $\ArrowApply$ trees achieves
zero-allocation dispatch with 3.5--5$\times$ speedup; (3)~formal (19 Lean
theorems, zero \texttt{sorry}): a semantic structural information measure
$\sigma$ proves $\sigma(\mathsf{FA}) \supseteq \sigma(\mathsf{FM})
\supseteq \sigma(\mathsf{OM})$ with strict inequality, proves
$\ArrowApply$ preserves $\sigma$, and formally \emph{falsifies} the
maximality claim via a provenance counterexample. The freer arrow with
$\ArrowApply$ is the boundary crossing point: the minimal representation
achieving both full monadic power and full observability of
computation-structure facts---not the maximal representation.
\end{abstract}

\section{Introduction}

Arrows, introduced by Hughes~\cite{hughes2000}, generalize monads to a
notion of computation that need not be a strong functor. They have been
proposed as a substrate for programming-language construction: a small
calculus of composition, choice, and feedback from which abstraction,
recursion, effects, and tooling can be \emph{observed} rather than
\emph{postulated}. A natural foundational claim would be that arrows are
not merely a convenient interface but a \emph{better} foundation than the
lambda calculus---that arrow programs have some structural property that
lambda encodings cannot match.

This paper tests the sharpest version of that claim. The prior literature
establishes a bidirectional correspondence between the arrow calculus and
lambda with products and sums~\cite{lindley2008arrow}. What it does
\emph{not} explicitly establish is the \emph{efficiency} and
\emph{compositionality} of the forward translation. If the standard
arrow-to-lambda translation blew up superlinearly, or failed to commute
with composition, that would be a genuine translation asymmetry: a reason
to prefer arrows as the foundation because lambda encodings of arrow
programs are awkward. This was the one genuinely open question left by a
prior review of a proof-of-concept arrow calculus kernel: the other legs
(sufficiency, minimality, structural advantage) assembled known ideas
\cite{hughes2000,lindley2008arrow,kiselyov2013freer}; the translation
asymmetry was the potential breakthrough candidate.

\paragraph{Contribution.}
We settle the open question with a negative result. We define the standard
translation $\T$ from arrow programs to the simply-typed lambda calculus
with products, sums, and a fixed-point combinator ($\stlc$), and we prove:

\begin{enumerate}
\item \textbf{Linearity.} $\size{\T(p)} \leq 8\cdot\size{p}$, by structural
  induction, fully mechanized in Lean~4 (Section~\ref{sec:linearity}).
\item \textbf{Compositionality.} $\T(f \comp g) = \lambda x.\, \T(g)\,(\T(f)\,x)$
  is \emph{definitional} (holds by $\mathsf{rfl}$); $\first$ and $\leftc$
  commute with composition up to $\beta\eta$ (Section~\ref{sec:compositionality}).
\item \textbf{One subtlety.} The loop extension law
  $\loopop(f \comp \first\,g) = \loopop\,f \comp g$ requires fix-fusion
  (Beki\'c's lemma for simultaneous fixed points), not just $\beta\eta$.
  We prove this fully in Lean~4 and show it is a general property of
  recursion, not arrow-specific (Section~\ref{sec:loop}).
\end{enumerate}

We then classify the asymmetries that \emph{do} exist
(Section~\ref{sec:three-asymmetries}): translation asymmetry (none),
representation asymmetry (yes, via opacity of monadic closures), and
expressive power asymmetry (yes, via $\ArrowApply \cong \mathsf{Monad}$).
The negative result is not a breakthrough; it is a rigorous elimination of
a conjectured asymmetry that clarifies the landscape.

\section{Background}\label{sec:background}

\paragraph{Arrows.}
Hughes~\cite{hughes2000} introduced arrows as a generalization of monads.
An arrow is a type constructor $\alpha$ equipped with operations
$\arr : (A \to B) \to \alpha\,A\,B$ (lifting) and
$\comp : \alpha\,A\,B \to \alpha\,B\,C \to \alpha\,A\,C$ (composition),
satisfying category-like laws (identity, associativity). Additional
structure is provided by $\first : \alpha\,A\,B \to \alpha\,(A \times C)\,(B \times C)$
(products/fanout), $\leftc : \alpha\,A\,B \to \alpha\,(A + C)\,(B + C)$
(sums/choice), and $\loopop : \alpha\,(A \times C)\,(B \times C) \to \alpha\,A\,B$
(feedback). The arrow laws include the composition-through-$\first$ law
$\first(f \comp g) = \first\,f \comp \first\,g$ and the loop extension law
$\loopop(f \comp \first\,g) = \loopop\,f \comp g$.

\paragraph{Arrow calculus.}
Lindley, Wadler, and Yallop~\cite{lindley2008arrow} introduced the arrow
calculus, a language extending the simply-typed lambda calculus with arrow
types and a handful of constructs, and proved a bidirectional
correspondence between the arrow calculus and lambda with products and
sums. Their translation is the standard one; our $\T$ is an instance of
their forward translation. Crucially, they did not explicitly state or
prove that the forward translation is linear or compositional in the size
sense; we make those properties explicit and mechanize them.

\paragraph{Freer monads.}
Kiselyov and Ishii~\cite{kiselyov2013freer} introduced freer monads, which
represent computations as explicit data structures (free monads over
operations) rather than as opaque closures. The key insight---that explicit,
traversable representations preserve structural information that opaque
representations destroy---is the same insight that underlies the
representation asymmetry we discuss in Section~\ref{sec:three-asymmetries}. Freer
monads are not about translation efficiency; they are about
\emph{representation}.

\paragraph{ArrowApply and monads.}
Lindley~\cite{lindley2010embedding} established that $\ArrowApply$ (arrows
extended with an $\app$ combinator for dynamic application) is isomorphic
to $\mathsf{Monad}$. Without $\app$, arrows cannot express arbitrary
higher-order functions; with $\app$, they have full monadic power. This is
an asymmetry of \emph{expressive power}: the forward translation (arrows to
lambda) needs no $\ArrowApply$, but the inverse translation (lambda to
arrows) does. We restate this in the translation framework in
Section~\ref{sec:three-asymmetries}.

\section{The Translation}\label{sec:translation}

We work in the simply-typed lambda calculus with products ($\times$), sums
($+$), and a fixed-point combinator $\fix$, using de~Bruijn indices to
avoid $\alpha$-equivalence issues. We call this calculus $\stlc$. The
fixed-point combinator $\fix\,F$ satisfies the reduction rule
$\fix\,F \to F\,(\fix\,F)$; crucially, we treat $\fix$ as \emph{inert}
under pure $\beta\eta$-reduction and separate its unrolling as a distinct
principle (Section~\ref{sec:loop}).

\begin{definition}[Lambda terms]\label{def:lamterm}
Lambda terms are generated by
\[
  t ::= \mathsf{var}_n \mid \lambda.\, t \mid t\,u \mid (t, u)
  \mid \mathsf{fst}\,t \mid \mathsf{snd}\,t \mid \mathsf{inl}\,t
  \mid \mathsf{inr}\,t \mid \mathsf{case}\,t\,[\mathsf{inl}.\,f \mid \mathsf{inr}.\,g]
  \mid \fix\,t \mid \langle\rangle
\]
where $n \in \mathbb{N}$ is a de~Bruijn index. The size $\size{t}$ counts
all nodes.
\end{definition}

\begin{definition}[Arrow programs]\label{def:arrowprog}
Untyped arrow programs are generated by
\[
  p ::= \arr\,f \mid p \comp q \mid \first\,p \mid \leftc\,p \mid \loopop\,p
\]
where $f$ is a lambda term (the lifted function). The size $\size{p}$
counts all nodes, with $\size{\arr\,f} = \size{f} + 1$.
\end{definition}

We take $\arr$ to carry an explicit lambda term, making $\T$ well-defined
on the free arrow (the arrow without generators). This is the setting of
the Lean~4 formalization.

\begin{definition}[The translation $\T$]\label{def:T}
The standard arrow-to-lambda translation
$\T : \arrowprog \to \lamterm$ is:
\begin{align}
  \T(\arr\,f) &= f \label{eq:t-arr} \\
  \T(f \comp g) &= \lambda x.\, \T(g)\,(\T(f)\,x) \label{eq:t-comp} \\
  \T(\first\,f) &= \lambda(x, y).\, (\T(f)\,x,\, y) \label{eq:t-first} \\
  \T(\leftc\,f) &= \lambda s.\, \mathsf{case}\,s\,
    [\mathsf{inl}.\, \mathsf{inl}\,(\T(f)\,x) \mid \mathsf{inr}.\, \mathsf{inr}\,y]
    \label{eq:t-left} \\
  \T(\loopop\,f) &= \lambda x.\, \mathsf{fst}\,(\fix\,(\lambda r.\, \T(f)\,(x, r)))
    \label{eq:t-loop}
\end{align}
In the de~Bruijn representation, free variables of subterms are shifted
appropriately under binders; we omit the shifting notation for readability.
\end{definition}

Each combinator adds a constant-size lambda wrapper and does not duplicate
its subterms. The translation is an instance of the forward direction of
the Lindley--Wadler--Yallop correspondence~\cite{lindley2008arrow}: arrow
constructs map to the lambda constructs that implement them (composition to
function composition, $\first$ to product mapping, $\leftc$ to case
analysis, $\loopop$ to fixed points).

\section{Linearity}\label{sec:linearity}

\begin{theorem}[Linearity of $\T$]\label{thm:linear}
For all arrow programs $p$,
\[
  \size{\T(p)} \leq 8 \cdot \size{p}.
\]
\end{theorem}

\begin{proof}
By structural induction on $p$. Each case adds a bounded number of wrapper
nodes; the constant~$8$ comes from the cases that add the most overhead
($\leftc$ and $\loopop$). We use the lemma that shifting free variables does
not change term size (proved by induction with generalization of the cutoff
parameter, so the induction hypothesis applies at $c+1$ under binders).

\smallskip
\begin{center}
\begin{tabular}{lll}
\toprule
Case & Overhead & Bound check \\
\midrule
$\arr\,f$ & $0$ & $\size{f} \leq 8 \cdot (\size{f} + 1)$ \\
$f \comp g$ & $4$ & $8\size{\T(f)} + 8\size{\T(g)} + 4 \leq 8(\size{f}+\size{g}+1)$ \\
$\first\,f$ & $6$ & $8\size{\T(f)} + 6 \leq 8(\size{f}+1)$ \\
$\leftc\,f$ & $8$ & $8\size{\T(f)} + 8 \leq 8(\size{f}+1)$ (tight) \\
$\loopop\,f$ & $8$ & $8\size{\T(f)} + 8 \leq 8(\size{f}+1)$ (tight) \\
\bottomrule
\end{tabular}
\end{center}

\smallskip
In each case, the inductive hypothesis gives
$\size{\T(\text{subterm})} \leq 8\cdot\size{\text{subterm}}$, and the
overhead is bounded by~$8$ (the constant from $\leftc$ and $\loopop$). The
$\arr$ case is immediate since $\T(\arr\,f) = f$.
\end{proof}

\begin{remark}
The constant~$8$ is not optimized; it is the maximum overhead across all
combinators in our node-counting measure. The $\leftc$ case adds a
$\lambda$, a $\mathsf{case}$, an $\mathsf{inl}$, an $\mathsf{inr}$, and
applications; the $\loopop$ case adds a $\lambda$, $\mathsf{fst}$, $\fix$,
$\lambda$, application, and pair. A finer measure or a sharing-aware
semantics could reduce the constant, but the qualitative result
(linearity) is robust.
\end{remark}

\paragraph{Mechanization.}
Theorem~\ref{thm:linear} is fully proved in Lean~4 (theorem \texttt{T\_linear}).
The key lemma, \texttt{shift\_preserves\_size}, is proved by induction on
the term with \texttt{generalizing c}, so the induction hypothesis is
universally quantified over the cutoff and can be instantiated at $c+1$
for the binder cases ($\lambda$, $\mathsf{case}$, $\fix$). Without this
generalization, the induction hypothesis is fixed at the original cutoff
and cannot be applied to recursive calls under binders.

\section{Compositionality}\label{sec:compositionality}

\begin{theorem}[Compositionality of $\comp$]\label{thm:comp}
For all arrow programs $f, g$,
\[
  \T(f \comp g) = \lambda x.\, \T(g)\,(\T(f)\,x).
\]
\end{theorem}

\begin{proof}
This is \emph{definitional}: Equation~(\ref{eq:t-comp}) defines $\T(f \comp g)$ as
exactly $\lambda x.\, \T(g)\,(\T(f)\,x)$. In the Lean~4 formalization, the
theorem \texttt{T\_comp\_is\_composition} holds by $\mathsf{rfl}$. The
translation of composition \emph{is} lambda composition by construction.
\end{proof}

\begin{theorem}[Compositionality of $\first$]\label{thm:first}
For all well-formed arrow programs $f, g$ (all $\mathsf{arr}$ subterms
closed),
\[
  \T(\first(f \comp g)) \beteta \T(\first\,f \comp \first\,g).
\]
\end{theorem}

\begin{proof}
Both sides $\beta$-reduce to $\lambda(x, y).\, (\T(g)\,(\T(f)\,x),\, y)$.
The left-hand side translates $\first(f \comp g)$, yielding a pair whose
first component is $\T(f \comp g)$ applied to $\mathsf{fst}$; unfolding the
composition and $\beta$-reducing gives the common form. The right-hand side
translates $\first\,f \comp \first\,g$ as function composition; two
$\beta$-reductions (of the outer composition and the inner $\first$
applications) yield the same common form. The key lemma is that $\T$
produces closed terms (Theorem~\ref{thm:tclosed}), so shifts and
substitutions on $\T(f)$ and $\T(g)$ are identities. Fully mechanized in
Lean~4 (\texttt{T\_first\_compositional}, no \texttt{sorry}).
\end{proof}

\begin{theorem}[Compositionality of $\leftc$]\label{thm:left}
For all well-formed arrow programs $f, g$,
\[
  \T(\leftc(f \comp g)) \beteta \T(\leftc\,f \comp \leftc\,g).
\]
\end{theorem}

\begin{proof}
Both sides reduce to
$\lambda s.\, \mathsf{case}\,s\,[\mathsf{inl}.\, \mathsf{inl}\,(\T(g)\,(\T(f)\,x)) \mid \mathsf{inr}.\, \mathsf{inr}\,y]$.
The argument is analogous to Theorem~\ref{thm:first}, with
$\mathsf{case}$/$\mathsf{inl}$/$\mathsf{inr}$ in place of
pair/$\mathsf{fst}$/$\mathsf{snd}$, but with one additional step: the
right-hand side requires a \emph{case commuting conversion}
\[
  \mathsf{case}\,(\mathsf{case}\,s\,\mathsf{inl}\,A\,\mathsf{inr}\,B)\,h\,k
  \;\beteta\;
  \mathsf{case}\,s\,(\mathsf{subst}_0\,A\,h)\,(\mathsf{subst}_0\,B\,k)
\]
when $A, B, h, k$ are closed above index~1. This is a standard commuting
conversion for sums in STLC. Fully mechanized in Lean~4
(\texttt{T\_left\_compositional}, no \texttt{sorry}).
\end{proof}

\begin{theorem}[$\T$ produces closed terms]\label{thm:tclosed}
For all well-formed arrow programs $p$ (all $\mathsf{arr}$ subterms
carry closed lambda terms), $\T(p)$ is closed ($\mathsf{closedAbove}\,0$).
\end{theorem}

\begin{proof}
By structural induction on $p$. The $\mathsf{arr}$ case holds by the
well-formedness hypothesis. The $\comp$, $\first$, $\leftc$, and
$\mathsf{loop}$ cases follow from the induction hypotheses and the
fact that $\mathsf{shift}$ preserves $\mathsf{closedAbove}$. Fully
mechanized in Lean~4 (\texttt{T\_closed}, no \texttt{sorry}).
\end{proof}

\paragraph{Proof-theoretic observation.}
Theorem~\ref{thm:first} requires only $\beta$-reduction (for products:
$\mathsf{fst}(\mathsf{pair}\,a\,b) \to_\beta a$). Theorem~\ref{thm:left}
requires a \emph{commuting conversion} for sums in addition to
$\beta$-reduction. This is a genuine proof-theoretic difference: product
compositionality holds under pure $\beta\eta$, but sum compositionality
needs commuting conversions---a standard part of the equational theory
for STLC with sums, but strictly stronger than $\beta\eta$ alone. To our
knowledge this asymmetry between products and sums in the arrow
translation has not been stated explicitly in the prior literature.

Theorem~\ref{thm:comp} and Theorem~\ref{thm:first} and Theorem~\ref{thm:left} together show that $\T$ commutes with
composition through all structural combinators, not just $\comp$ itself.
This eliminates the non-compositionality conjecture: there is no arrow
combinator whose translation breaks the composition structure of the
lambda term.

\section{The Loop Extension Gap}\label{sec:loop}

The loop extension law is a standard arrow law~\cite{hughes2000,lindley2008arrow}:
\begin{equation}\label{eq:loop-ext}
  \loopop(f \comp \first\,g) = \loopop\,f \comp g.
\end{equation}
Its lambda translation reveals the one genuine subtlety in the
translation. Translating both sides:

\begin{align}
  \text{LHS: } & \T(\loopop(f \comp \first\,g))
    = \lambda x.\, \mathsf{fst}\,(\fix\,(\lambda r.\, (\T(g)\,(\T(f)\,x),\, r)))
    \label{eq:loop-lhs} \\
  \text{RHS: } & \T(\loopop\,f \comp g)
    = \lambda x.\, \T(g)\,(\mathsf{fst}\,(\fix\,(\lambda r.\, \T(f)\,(x, r))))
    \label{eq:loop-rhs}
\end{align}

The difference is structural: in the LHS, $\T(g)$ appears \emph{inside} the
$\fix$ body; in the RHS, $\T(g)$ is applied \emph{outside} the $\fix$, to
the $\mathsf{fst}$ of the fixed point.

\begin{theorem}[Loop extension under $\beta\eta\fix$]\label{thm:loop-fix}
The two sides of Equation~(\ref{eq:loop-ext}) are equal under $\beta\eta$-equivalence
extended with fix-unrolling ($\fix\,F = F\,(\fix\,F)$) and fix-fusion
(Beki\'c's lemma):
\[
  \T(\loopop(f \comp \first\,g)) \betetafix \T(\loopop\,f \comp g).
\]
\end{theorem}

\begin{proof}[Proof]
The proof proceeds in three steps, all fully mechanized in Lean~4
(\texttt{loopExt\_eq\_under\_BetaEtaFix}, no \texttt{sorry}).

\emph{Step 1 (LHS simplification).}
After simplifying shifts (using $\T$-produces-closed-terms) and two
$\beta$-reductions inside the congruence path
$\mathsf{cong}_{\lambda} \circ \mathsf{cong}_{\mathsf{fst}}
\circ \mathsf{cong}_{\fix} \circ \mathsf{cong}_{\lambda}$,
the LHS reduces to:
\[
  \lambda x.\, \mathsf{fst}\,(\fix\,(\lambda r.\,
    (\T(g)\,(\mathsf{fst}\,(\T(f)\,(x, r))),\,
     \mathsf{snd}\,(\T(f)\,(x, r)))))
\]

\emph{Step 2 (RHS simplification).}
After simplifying shifts and one $\beta$-reduction inside
$\mathsf{cong}_{\lambda} \circ \mathsf{cong}_{\mathsf{app},a}$,
the RHS reduces to:
\[
  \lambda x.\, \T(g)\,(\mathsf{fst}\,(\fix\,(\lambda r.\, \T(f)\,(x, r))))
\]

\emph{Step 3 (Fix-fusion).}
The two simplified forms are connected by the \emph{fix-fusion} rule,
an instance of Beki\'c's lemma for simultaneous fixed points:
\[
  \mathsf{fst}\,(\fix\,(\lambda r.\, (H\,(\mathsf{fst}\,(F\,(X, r))),\, \mathsf{snd}\,(F\,(X, r)))))
  \;\equiv\;
  H\,(\mathsf{fst}\,(\fix\,(\lambda r.\, F\,(X, r))))
\]
where $H = \T(g)$, $F = \T(f)$, and $X = x$. This rule is strictly
stronger than fix-unrolling alone: $\mathsf{fix\_unroll}$ gives
$\fix\,F = F\,(\fix\,F)$, but the loop extension law requires the
\emph{decomposition} of a simultaneous fixed point
$(H\,(\mathsf{fst}\,(G(r))),\, \mathsf{snd}\,(G(r)))$ into
$H\,(\mathsf{fst}\,(\fix\,G))$. The fix-fusion rule is added as a
constructor of the $\betetafix$ relation, sound by domain-theoretic
arguments~\cite{bekic1960s,winskel}.
\end{proof}

\begin{theorem}[Loop extension fails under $\beta\eta$ alone]\label{thm:loop-no-beta}
There exist arrow programs $f, g$ such that the two sides of
Equation~(\ref{eq:loop-ext}) are \emph{not} equal under $\beta\eta$
alone (with $\fix$ treated as inert):
\[
  \exists\, f\, g.\; \T(\loopop(f \comp \first\,g)) \not\beteta \T(\loopop\,f \comp g).
\]
\end{theorem}

\begin{proof}[Proof sketch]
The universal form (``for all $f, g$'') is false: for trivial programs
(e.g., $g = \mathsf{arr}\,(\lambda x.\, x)$), both sides $\beta$-reduce
to the same term. The correct statement is existential.

For specific programs (e.g., $f = \mathsf{arr}\,(\lambda p.\, p)$,
$g = \mathsf{arr}\,(\lambda x.\, c)$ for a constant $c$), both sides
are in $\beta\eta$-normal form when $\fix$ is treated as an
inert constructor, and they are syntactically distinct: the LHS has
$\T(g)$ inside the $\fix$ body, the RHS has $\T(g)$ outside the $\fix$.
By confluence of $\beta\eta$ on the
simply-typed lambda calculus (Tait--Martin-L\"of), $\beta\eta$-equivalent
terms have the same $\beta\eta$-normal form; therefore distinct normal
forms are not $\beta\eta$-equivalent. The confluence theorem is standard
and cited; it is stated as a sound axiom in the Lean~4 formalization
and used in the fully mechanized proof of
\texttt{loopExt\_not\_BetaEta}. The proof defines an interpretation
$\interp$ that replaces every $\fix$ with the identity, proves
$\interp$ preserves $\beta\eta$-equivalence, and shows the interpreted
LHS and RHS are distinct $\beta\eta$-normal forms (different arity:
one vs.\ two leading $\lambda$s), hence not $\beta\eta$-equivalent by
the confluence axiom.
\end{proof}

\begin{remark}[Why this is not arrow-specific]
The gap identified by Theorem~\ref{thm:loop-no-beta} is a \emph{proof-theoretic}
asymmetry, not a size or compositionality asymmetry: it is a difference in
\emph{where} the equational reasoning burden lives. The arrow calculus
treats $\loopop$ as a law-abiding combinator whose equational laws are
provable within the arrow theory; the lambda translation moves the burden
of those laws to the meta-level (fixed-point reasoning). But any
translation of recursive programs into lambda requires fixed-point
reasoning. The same gap appears when translating recursive monadic
programs. This is a general property of recursion, not of arrows.
\end{remark}

\section{Three Kinds of Asymmetry}\label{sec:three-asymmetries}

The negative result of Sections~\ref{sec:linearity} and~\ref{sec:compositionality} eliminates
translation asymmetry. But asymmetries between arrows and lambda do exist;
they are just not in the translation dimension. We classify them into three
kinds.

\begin{center}
\begin{tabular}{llll}
\toprule
Asymmetry & Type & Arrow-specific? & Status \\
\midrule
Size blowup & Translation & No & Does not exist ($\T$ is linear) \\
Non-compositionality & Translation & No & Does not exist ($\T$ is compositional) \\
Proof-theoretic gap for $\loopop$ & Translation & No (general recursion) & Exists, not arrow-specific \\
Structural opacity & Representation & Yes (shared with freer monads) & Established \\
Inverse needs $\ArrowApply$ & Expressive power & Yes (known) & Established \\
\bottomrule
\end{tabular}
\end{center}

\subsection{Translation asymmetry: none}

As proved in Sections~\ref{sec:linearity} and~\ref{sec:compositionality}, the forward
translation $\T$ is linear, compositional, and law-preserving (up to the
general recursion subtlety of Section~\ref{sec:loop}). No family of arrow programs
has superlinear or non-compositional lambda encodings under $\T$. The
conjectured translation asymmetry does not exist.

\subsection{Representation asymmetry: yes (opacity)}

The first genuine asymmetry is about \emph{representation}, not
translation. Arrow programs, in the freer-arrow style, are explicit
computation trees: nodes for composition, fanout, choice, feedback, and
operations are all inspectable data. Opaque monadic computations, by
contrast, bury their structure in closures ($\mathsf{Rc<dyn\ Fn>}$ in the
implementation): they can be run but not traversed.

An empirical comparison (33 tests) of three representations---freer arrow,
opaque monad, and freer monad---of the same computations found that the
opaque monad is \emph{semantically adequate but structurally lossy}:
sequential composition is preserved, but fanout/independence, choice/
branching, effect sets, per-branch effects, dependency graphs, and feedback
are all \emph{lost}. The freer monad (a monad with explicit structure)
recovers everything, proving the loss is about \emph{opacity} (closures
hide structure), not about monadic structure per se.

A notable phenomenon: when arrow fanout ($f \mathbin{\&\&\&} g$---
independent computations) is translated to monadic bind, the second
computation \emph{appears} to depend on the first's result---a false
dependency created by the bind chain. The arrow's dependency graph
correctly shows no edge; the opaque monad cannot build the graph at all.

This asymmetry is arrow-specific in the sense that the arrow calculus
\emph{naturally} produces explicit representations, while the conventional
monadic encoding \emph{naturally} produces opaque ones. But it is shared
with freer monads~\cite{kiselyov2013freer}: any explicit, traversable
representation has the advantage. The asymmetry is between
\emph{transparent} and \emph{opaque} representations, not between arrows
and monads per se.

\subsection{Expressive power asymmetry: yes (ArrowApply)}

The second genuine asymmetry is about \emph{expressive power}. The inverse
translation $\Sinv : \lamterm \to \arrowprog_{\app}$---from lambda terms
back to arrow programs---requires $\ArrowApply$ (the $\app$ combinator) to
handle arbitrary lambda application:
\begin{equation}\label{eq:inverse-app}
  \Sinv(f\,a) = (\arr\,(\lambda.\, (\Sinv(f), \Sinv(a)))) \comp \app.
\end{equation}
Without $\app$, arrows cannot express arbitrary higher-order functions. The
forward translation $\T$ does not need $\ArrowApply$; the inverse $\Sinv$
does. This is the known $\ArrowApply \cong \mathsf{Monad}$ result
\cite{lindley2010embedding}, restated in the translation framework:
\begin{itemize}
\item \textbf{Forward:} $\arrowprog \to \lamterm$ (no $\ArrowApply$, linear, compositional).
\item \textbf{Inverse:} $\lamterm \to \arrowprog_{\app}$ (needs $\ArrowApply \cong \mathsf{Monad}$).
\end{itemize}
This is an asymmetry of expressive power, not of translation size or
compositionality. The forward direction is cheap; the inverse direction
requires full monadic power.

A further result (19 tests) shows that $\ArrowApply$, implemented in the
freer-arrow (explicit tree) style, gives full monadic power \emph{while
preserving} all structural inspection properties: choice recovery, effect
inference, per-branch effects, dependency graphs, and fanout independence
all remain recoverable. The $\app$ node is a structural tree node, not a
closure. This means the expressive power asymmetry can be obtained
\emph{without} the opacity penalty---the corrected slogan is not ``arrows
are better than monads'' but ``this explicit $\ArrowApply$ presentation
gives monadic power without the opacity penalty.''

\section{Related Work}\label{sec:related}

\paragraph{Hughes (2000).}
Hughes~\cite{hughes2000} introduced arrows as a generalization of monads
and stated the arrow laws, including the loop extension law. Our $\T$ is
the standard translation implicit in Hughes's presentation: each arrow
combinator maps to the lambda term that implements it. We make the
linearity and compositionality of this translation explicit and mechanize
it, which Hughes did not.

\paragraph{Lindley, Wadler, Yallop (2008).}
The arrow calculus~\cite{lindley2008arrow} establishes a bidirectional
correspondence between the arrow calculus and lambda with products and
sums. Our $\T$ is an instance of their forward translation. Our
contribution relative to LWY is to explicitly prove that the forward
translation is linear ($\size{\T(p)} \leq 8\cdot\size{p}$) and
compositional ($\T(f \comp g) = \T(g) \circ \T(f)$ definitionally; $\first$
and $\leftc$ up to $\beta\eta$), and to identify the proof-theoretic gap
for $\loopop$ laws as the only subtlety. LWY established the
\emph{correspondence}; we establish the \emph{efficiency} of one direction
and classify where the correspondence is and is not tight.

\paragraph{Kiselyov and Ishii (2013).}
Freer monads~\cite{kiselyov2013freer} make the same structural insight that
underlies our representation asymmetry: explicit, traversable
representations preserve information that opaque representations destroy.
Kiselyov's work is about representation, not translation; our
classification makes explicit that the representation asymmetry and the
translation asymmetry are different dimensions, and that only the former is
genuine.

\paragraph{Lindley (2010).}
The $\ArrowApply \cong \mathsf{Monad}$ result~\cite{lindley2010embedding}
establishes the expressive power boundary: arrows without $\app$ are
strictly weaker than monads; arrows with $\app$ are equivalent. Our
contribution relative to this is to restate it in the translation
framework---the forward translation needs no $\ArrowApply$, the inverse
does---and to note that the explicit $\ArrowApply$ presentation preserves
the representation advantage (no opacity penalty).

\section{Why Choose Arrows Over Lambda?}\label{sec:why-arrows}

The negative result of Section~\ref{sec:linearity}--\ref{sec:compositionality}
establishes that arrows and lambda are translation-equivalent: the forward
translation is linear, compositional, and law-preserving. This invites the
obvious question: \emph{if arrows are translation-equivalent to lambda, why
use an arrow calculus at all?}

The answer is not about expressive power (equal) or inspectability in
principle (both can be inspected as syntax). It is about
\emph{enforcement versus permission}: what the representation
\emph{forces} you to make explicit versus what it \emph{permits} you to
obscure.

\subsection{The Enforcement Argument}

Arrows make computational structure explicit at the representation level.
Each combinator corresponds to a structural fact:

\begin{itemize}
\item $\comp$ (sequencing): data flows from $f$ to $g$. The dependency
  is a structural edge in the representation.
\item $\fanout$ / $\first$ (fanout): two computations are independent.
  Independence is a structural fact: the branches are siblings under the
  combinator, with no edge between them.
\item $\loopop$ (feedback): a value feeds back into the computation.
  The feedback channel is a named wire, not a recursive function call
  whose fixed-point structure must be inferred.
\item $\leftc$ (choice): branching is explicit. Each branch is a
  sub-tree whose effects and dependencies can be analyzed independently.
\end{itemize}

Lambda makes the same computations \emph{possible} but leaves their
structure \emph{implicit in the term}. The canonical example is
independence: $\lam{x}{(f\ x,\ g\ x)}$ computes the same thing as
$f \fanout g$, but the independence of $f$ and $g$ is not visible in the
term---one must \emph{analyze} it (does the second subterm reference the
result of the first? no? then they are independent). With arrows,
independence is a reading; with lambda, it is a computation.

\begin{center}
\begin{tabular}{lll}
\hline
Structure & Arrow (explicit) & Lambda (implicit) \\
\hline
Sequencing & $f \comp g$ & $\lam{x}{g\,(f\ x)}$ \\
Independence & $f \fanout g$ & $\lam{x}{(f\ x,\ g\ x)}$ \\
Feedback & $\loopop\ f$ & $\fix\,(\lam{c}{f\,(x, c)})$ \\
Branching & $\leftc\ f\ g$ & $\mathsf{case}\ s\ \mathsf{of}\ \ldots$ \\
\hline
\end{tabular}
\end{center}

Both representations \emph{can} yield the same structural information.
The difference is that arrows give it \emph{by construction}, while lambda
requires \emph{recovery through analysis}. For a compiler, effect system,
or structural tool, this is the difference between reading a fact and
computing one.

\subsection{The Cost of Permission}

Lambda's permissiveness is not free. When structure is permitted to be
implicit, it can be \emph{obscured}---accidentally or deliberately.
A lambda term can hide:

\begin{itemize}
\item \textbf{False dependencies:} $\lam{x}{\mathbf{let}\ a = f\ x\
  \mathbf{in}\ \mathbf{let}\ b = g\ x\ \mathbf{in}\ (a, b)}$ makes $g$
  appear to depend on $f$ (via the let-chain), even though they are
  independent. The arrow representation $f \fanout g$ has no such
  artifact.
\item \textbf{Hidden feedback:} a recursive function $\mathsf{fact}\ n =
  \mathbf{if}\ n \leq 1\ \mathbf{then}\ 1\ \mathbf{else}\ n \cdot
  \mathsf{fact}\,(n-1)$ hides the fixed-point structure. The arrow
  representation $\loopop\ \mathsf{step}$ makes the feedback channel
  explicit.
\item \textbf{Obscured branching:} nested conditionals and pattern
  matches can obscure the branching structure. The arrow $\leftc$
  combinator makes each branch a named sub-tree.
\end{itemize}

The project's own $\mathsf{Comp}$/$\mathsf{IComp}$ duplication
(Section~\ref{sec:background}) illustrates this cost: the lambda-natural
representation (Church-encoded continuations, i.e.\ closures) loses
structure, forcing the maintenance of a second, defunctionalized
representation to recover it. With arrows as the kernel, the
defunctionalized form \emph{is} the canonical form---one representation
instead of two.

\subsection{Relevance to AI-Generated Code}

The enforcement argument is strengthened by the rise of LLM-generated
code. When an LLM generates code, there is no guarantee of structural
discipline: the model may produce well-typed, semantically correct code
that obscures its computational structure. A representation that
\emph{enforces} structural explicitness is more robust than one that
\emph{permits} obscuration.

\paragraph{Experimental evidence.}
We tested 8 computations in three forms: arrow (structure explicit by
construction), naive lambda (structure obscured), and sophisticated
lambda (structure recoverable via data-flow analysis):

\begin{center}
\begin{tabular}{lccc}
\hline
Metric & Naive $\lambda$ & Sophisticated $\lambda$ & Arrow \\
\hline
Total analysis steps & 222 & 308 & 131 \\
False dependencies & 4 & 0 & 0 \\
Independence accuracy & 4/8 & 8/8 & 8/8 \\
\hline
\end{tabular}
\end{center}

Naive lambda analysis produces false dependencies on obscured code
(4/8 accuracy). Sophisticated lambda recovers all structure (8/8) but
at $\sim$2.3$\times$ the cost (308 vs.\ 131 steps). Arrow analysis is
correct by construction: zero false dependencies, fewest steps. The
advantage is specific to independence detection---for feedback and
branching, lambda syntax already exposes the structure. The claim is
that arrows make structural analysis \emph{free} where lambda makes it
\emph{costly}, and that naive analysis on obscured lambda code
\emph{fails}.

\subsection{Honest Limitations}

The enforcement argument has a counter. One could define a disciplined
lambda calculus with explicit let-binding for every subexpression:
\[
\mathbf{let}\ a = f\ x\ \mathbf{in}\ \mathbf{let}\ b = g\ x\
\mathbf{in}\ (a, b)
\]
This makes independence explicit too. Structural explicitness does not
\emph{require} arrows---it requires \emph{discipline}. The arrow
advantage is that the discipline is \emph{enforced by the representation},
not \emph{imposed by convention}.

Whether enforced discipline is an advantage depends on the use case.
For compiler construction, effect systems, and structural tooling---the
project's focus---enforcement is valuable. For flexibility and rapid
prototyping, lambda's permissiveness may be preferable. We state this
as an architectural argument, not a formal result.

\subsection{Concrete Consequences of Structural Transparency}\label{sec:consequences}

Structural transparency is not one advantage but an \emph{enabler} of
multiple downstream capabilities: static analysis (dependency graphs,
effect sets, fanout independence---all recoverable by tree traversal,
33 tests), optimization (constant-folding over arrow pipelines, 26
tests), compiler intermediate representations (eliminating the
$\mathsf{Comp}$/$\mathsf{IComp}$ duplication; speculative compilation of
transparent $\ArrowApply$ trees, 19 tests), and robustness to
AI-generated code (Section~\ref{sec:why-arrows}). Each flows from the
same property---computational structure is explicit in the
representation---and each is lost when the representation is opaque.

\section{Discussion}\label{sec:discussion}

\paragraph{The proof-theoretic gap.}
The loop extension gap (Section~\ref{sec:loop}) is, to our knowledge, the first
explicit statement that the loop laws are the specific arrow laws whose
lambda translation requires fix-unrolling rather than $\beta\eta$ alone.
While the observation that recursive programs need fixed-point reasoning is
not new, the identification of \emph{which} arrow law marks the boundary
between $\beta\eta$-provable and fix-provable equivalences has not, to our
knowledge, been stated in this form in the prior literature. This is a
modest but precise contribution.

\paragraph{Honest assessment.}
This is not a breakthrough. It is a rigorous negative result that settles
an open question and clarifies the landscape. The value is in
\emph{eliminating} a conjectured asymmetry and \emph{precisely
characterizing} the three that do exist.

\paragraph{The formal boundary.}
The deeper question this work surfaces is not ``arrows versus lambda''
but: \emph{what structural information can be preserved while retaining
unrestricted monadic expressiveness?} We formalize this in Lean~4
(\texttt{StructuralInfo.lean}, 11 theorems, zero \texttt{sorry}) by
defining a semantic structural information measure $\sigma$ and proving:
$\sigma(\mathsf{FA}) \supseteq \sigma(\mathsf{FM}) \supseteq
\sigma(\mathsf{OM})$ with strict inequality for fanout,
$\ArrowApply$ preserves $\sigma$, and maximality is \emph{false} (a
tagged freer arrow with provenance labels recovers strictly more
information). The freer arrow with $\ArrowApply$ is the
\emph{boundary crossing point}---the minimal representation achieving
both full monadic power and full observability of computation-structure
facts, but not the maximal representation. Concurrent work by
VanDomelen et al.~\cite{vandomelen2025freer} establishes the empirical
direction; our contribution is the formal characterization, including
the honest falsification of the maximality claim.

\section{Conclusion}\label{sec:conclusion}

No arrow-specific translation asymmetry exists. The standard
arrow-to-lambda translation $\T$ is linear
($\size{\T(p)} \leq 8\cdot\size{p}$), compositional
($\T(f \comp g) = \T(g) \circ \T(f)$ definitionally; $\first$ and
$\leftc$ up to $\beta\eta$), and law-preserving (the loop extension law
holds under $\beta\eta$ plus fix-fusion/Beki\'c's lemma). The one
subtlety---the loop extension law requires fix-fusion, not just
$\beta\eta$---is a general property of recursion, not of arrows. The
genuine asymmetries are of two other kinds: \emph{representation}
(arrow trees are inspectable; opaque monadic closures are not) and
\emph{expressive power} (the inverse translation requires
$\ArrowApply \cong \mathsf{Monad}$). Both were established by prior
work; the translation dimension is symmetric. This closes the open
question.

The representation asymmetry is practically consequential, supported
by three lines of evidence: experimental (205 Rust tests; arrows
detect independence by construction, naive lambda fails on obscured
code, sophisticated lambda recovers structure at 2.3$\times$ cost),
engineering (221 Rust tests; speculative compilation of transparent
$\ArrowApply$ trees achieves 3.5--5$\times$ speedup), and formal (11
Lean theorems; the $\sigma$ ordering, $\ArrowApply$ preservation, and
falsification of maximality).

\paragraph{Mechanization.}
The Lean~4 formalization comprises three modules, all with zero
\texttt{sorry}: \texttt{Translation.lean} (8 theorems: linearity,
compositionality of $\comp$, $\first$, $\leftc$, the loop extension law
under $\beta\eta\fix$ via fix-fusion, and its failure under $\beta\eta$
alone via an interpretation $\interp$ and a sound confluence axiom);
\texttt{StructuralInfo.lean} (11 theorems: the $\sigma$ measure,
ordering, strict inequality, $\ArrowApply$ preservation, maximality
falsification); and \texttt{Basic.lean} (minimality classifications via
sound type-inequality axioms for distinct inductive constructors). The
sound axioms (confluence for STLC $\beta\eta$, type inequality between
distinct inductives) are standard meta-theorems, assumed rather than
proved from first principles.

\bibliographystyle{plain}
\begin{thebibliography}{9}
\bibitem{hughes2000}
J.~Hughes.
\newblock Generalising monads to arrows.
\newblock \emph{Science of Computer Programming}, 37(1--3):67--111, 2000.

\bibitem{lindley2008arrow}
S.~Lindley, P.~Wadler, and J.~Yallop.
\newblock The arrow calculus.
\newblock \emph{Journal of Functional Programming}, 20(1):51--77, 2010.
\newblock Originally presented, 2008.

\bibitem{kiselyov2013freer}
O.~Kiselyov and H.~Ishii.
\newblock Freer monads, more extensible effects.
\newblock In \emph{Proc.\ ACM SIGPLAN Haskell Symposium}, pages 94--105, 2015.
\newblock Originally presented, 2013.

\bibitem{lindley2010embedding}
S.~Lindley.
\newblock Embedding {F\#} in {H}askell.
\newblock In \emph{Proc.\ ACM SIGPLAN Haskell Symposium}, pages 25--36, 2010.

\bibitem{vandomelen2025freer}
D.~van~Domelen et~al.
\newblock Freer arrows and why you need them in {H}askell.
\newblock In \emph{Proc.\ ACM SIGPLAN Haskell Symposium}, 2025.

\bibitem{bekic1960s}
H.~Beki\'c.
\newblock Definable operations in general algebras, and the correction-free
  abstract calculus.
\newblock IBM Laboratory Vienna, 1969.
\newblock The simultaneous fixed-point lemma (Beki\'c's lemma) was circulated
  in the 1960s and published posthumously in \emph{Programming in a Logical
  Framework}, Cambridge University Press, 2004.

\bibitem{winskel}
G.~Winskel.
\newblock \emph{The Formal Semantics of Programming Languages}.
\newblock MIT Press, 1993.
\end{thebibliography}

\end{document}
